\section{Study 1 details}\label{app:compiled}

\paragraph{Programs and site.}
Lookup programs aggregate the tokens at positions 0 and 1 into categorical variables and map them through the
table with a categorical MLP; additive programs aggregate the tokens categorically and map each to its numerical
value before a numerical MLP adds them, because the pinned Tracr compiler does not aggregate non-binary numerical
values. Compiler weights are used unchanged. Tables were drawn by rejection sampling with seeds 10000--10039
(lookup) and 20000--20039 (additive); no table repeats within or across families.

\paragraph{Labels.}
Let $Q$ be the rotation of the edited coordinates and let $C_A$ and $C_B$ select the unrotated carrier coordinates
of $a$ and $b$. An alignment is valid if, for every difference $d=z_s-z_r$ between natural states,
$dP_AQ^{\top}C_A=dQ^{\top}C_A$ and $dP_AQ^{\top}C_B=0$, and symmetrically for $P_B$, within $10^{-4}$, and
$\|P_AP_B-P_BP_A\|_{\max}<10^{-8}$. In words, an $A$-edit moves the $a$-carrier by exactly the donor's difference
and leaves the $b$-carrier alone. Because the final MLP reads only the carriers among the edited coordinates,
valid alignments implement every single and joint interchange correctly. The labels never use a test's output.

\paragraph{Alignments.}
Mean-contrast alignments take the singular vectors (singular value above $10^{-6}$) of the centred conditional means
of the site states for each value of the variable. The DAS alignment uses an orthogonal matrix parameterised by a
Cayley transform, initialised from the restricted mean contrasts plus a seeded random completion, with rank 2 per
variable for lookup and 1 for additive. It minimises the squared error of the output scores over the 144 single
interchanges among the 12 restricted inputs, with Adam (learning rate 0.02) for at most 200 steps, keeping the
lowest-loss iterate and stopping at a loss of $10^{-12}$, which keeps an exact solution from drifting at round-off
level. Every final DAS alignment reproduces all its training counterfactuals. Random alignments have the DAS ranks.

\paragraph{Tests.}
Each test draws its queries from a fixed seeded stream, the same for all alignments of a program; singles alternate
between $A$ and $B$; the joint, raw-distance, transport and natural tests share one stream of $(r,s,t)$ triples in
which both variables change. A test accepts an alignment if every sampled query passes. All allocated queries were
executed. Budgets count three calls per single query, four per joint and raw-distance query, six per transport query
and one per natural execution.

\paragraph{Development, freeze and amendment.}
Development used four programs per family. It already showed that broader single panels and the strong controls
rejected the restricted alignments as often as joint panels did; the protocol was then frozen and the final
programs generated. After the final run an independent re-execution found \CmpNearTies{} queries in which the
saved category and the recomputed one differ; in all of them the two classes' scores differ by at most
\CmpTieGap{}. A post-freeze amendment therefore treats a score gap within $10^{-4}$ as undecided. Under that rule
every rejection by the broader panels stands, and one rejection by restricted singles at 120 calls (a random
alignment) becomes undecided. The census in the main text uses the same rule.

\paragraph{Verification and replay.}
A separate verifier, which shares no model or panel code with the study, recomputes the networks head by head from
the saved weights, re-derives every label, re-executes all \CmpQueries{} saved queries with the margin rule, and
checks joint composition over all 19,683 triples for every alignment of the first program in each family. A fresh
environment with the pinned packages recompiled one program per family with bitwise-identical parameters and
reproduced all \CmpReplayRecords{} of their query records exactly, reusing the saved alignments.

\section{Study 2 details}\label{app:trained}
\paragraph{Tasks and networks.}
Tables were drawn with the generator of Study~1 from seeds 30000--30024 (lookup) and 40000--40024 (additive), which
no earlier analysis had used. The output is a class over the distinct values of the table. Transformers use two heads
of width $d/2$, LayerNorm before the attention and MLP sublayers and before the unembedding, and an MLP of width $4d$;
the attention mask is causal, so the last position reads the whole sequence. Networks were trained in double
precision on the full batch of 27 inputs with AdamW without weight decay, for 3,000 steps, with learning rate
$3\times10^{-3}$ for transformers and $10^{-2}$ for MLPs. The initialisation seed was ten times the table seed plus
the attempt number; up to three attempts were allowed, and the first network that classified all 27 inputs with a
logit margin of at least 1 was used. All \TrnCompetent{} networks passed, \TrnFirstAttempt{} at the first attempt.

\paragraph{Alignments.}
DAS parameterises an orthogonal matrix by a Cayley transform of a skew-symmetric matrix, initialised from the
orthonormalised conditional-mean contrasts of $A$ and $B$ completed with seeded random directions, and minimises the
cross-entropy of the counterfactual class over its training interchanges with Adam (learning rate 0.02, 1,500 steps),
keeping the lowest-loss iterate. The mean-contrast alignment takes, for each variable, the singular vectors of the
centred conditional means with singular values above $10^{-6}$ of the largest, and makes the two sets orthogonal by
symmetric orthogonalisation. Off-natural distances are Euclidean distances between an edited state and the natural
state of the input it simulates, divided by the root-mean-square distance between distinct natural states of the
same network.

\paragraph{Development and freeze.}
Before development, a small exploratory run trained twelve MLPs on lookup tables of Study~1 and found no joint-only
failures among their alignments; none of its networks or tables is reused.
Development used tables with seeds 50--53 in each family and fixed every setting above. It showed that alignments in
trained networks are approximate, that weight decay made them worse, and that longer DAS training did not improve
them; the primary quantities were therefore defined for approximate alignments. The development results, including
the sign of the primary contrast, are recorded in the protocol, which was frozen with the code and the final seeds
before any final network was trained.

\paragraph{Verification.}
An independent implementation in NumPy, sharing no code with the study, rebuilt every network from its saved weights,
confirmed that all of them classify the 27 inputs correctly, recomputed all \TrnVerifiedDecisions{} census decisions
(\TrnVerifiedDiffer{} differed), and recounted the per-alignment quantities of the analysis.


\section{Study 3 details}\label{app:multisite}
\paragraph{Tasks, networks and sites.}
Tables were drawn with the generator of Studies~1 and~2 from seeds 50000--50024 (lookup) and 60000--60024 (additive),
which no earlier analysis had used, and networks were trained by Study~2's code with its settings, including up to
three initialisations to classify all 27 inputs with a logit margin of at least 1 (Appendix~\ref{app:trained}).
Of the \MsNetworks{} networks, \MsCompetent{} reached this margin, \MsFirstAttempt{} at the first attempt. For the MLPs, the natural states at $A$'s site are the first hidden layer after its ReLU and those at $B$'s
site the second, and the output is a linear function of the latter. For the transformers, the natural state at $A$'s
site is the sum of the token and position embeddings of $a$ at position~1; an $A$-edit replaces it, the attention
layer is recomputed with the edited position, and the last position's residual stream after attention is $B$'s site.
Applying the propagation to the unedited states reproduces the natural states at $B$'s site exactly for every network.

\paragraph{Alignments.}
Each DAS projector is fitted at its own site with Study~2's settings (Cayley rotation initialised from the variable's
conditional-mean contrasts at that site with seeded random completion; Adam, learning rate 0.02, 1,500 steps; the
lowest-loss iterate), on the 486 single interchanges of its variable in which the variable changes. The mean-contrast
alignment projects onto the conditional-mean contrasts of $A$ at $A$'s site and of $B$ at $B$'s site; the random one
onto seeded random subspaces of rank 2. Off-target distances are divided by the root-mean-square distance between the
natural states of distinct inputs at the site concerned.

\paragraph{Development, predictions and freeze.}
The two predictions were written, with the rest of the protocol, before any network of this study existed, and the
time and SHA-256 of the protocol and of its predictions were recorded then. Development used tables with seeds 70--73
in each family. All \MsDevNetworks{} development networks reached the margin; \MsDevPerfect{} of the \MsDevAlignments{} alignments passed every single
query and none of them failed a joint query; the detection contrast at 120 calls was \MsDevContrast. DAS at the MLP's
second hidden layer converged within about 100 steps, and longer training or a larger learning rate gave the same best
losses on four networks; no setting was changed. The development results were appended to the protocol, which was
frozen with the code, the final seeds and the shared Study~2 module that the code imports before any final network was trained.

\paragraph{When the margin scores were fixed.}
On 6 October 2026 (UTC), the Study~2 script that defines the unsigned score $\rho$, the signed score $\sigma$ and the
rule $\sigma>1$ was written at 07:23 and first run at 07:24; the counts of queries flagged by the rule were added at
07:25. Study~3's predictions were recorded at \MsPredictionsTime, its code, which computes both one-term scores, was
frozen at 07:55, and its final networks finished training at 08:34. The alignment-level application of the rule to
Study~3 and the breakdown by architecture were coded after that, without changing the rule.

\paragraph{Verification.}
An independent implementation in NumPy, sharing no code with the study, rebuilt every network and both sites from the
saved weights, recomputed all \MsVerifiedDecisions{} census decisions (\MsVerifiedDiffer{} differed), and recounted
the per-alignment failures, joint-panel failures and unexplained joint failures of the analysis.


\section{Study 4: how the protocol changed, and when}\label{app:gpt2-history}

The protocol was approved as written, with \NTasksConf{} confirmation tasks and the sign rule of
\cref{sec:gpt2-design} as the primary pass rule. Its predictions, configuration and prompt lists were hashed before any
model was run. Every later change was recorded before the data it concerns.

\paragraph{Amendment 1: the parity gate (after the pilot, before development).} The pilot compared our GPT-2 code with
the reference implementation of the transformers library on two development draws. In float32 on the GPU their margins
differed by up to \PilotMarginGap, against a prespecified limit of $10^{-4}$; in float64 on the CPU they agreed to
\ParityFloatSixtyFour{} in every logit, so the gap was rounding. The gate became float64 agreement within $10^{-8}$ and
float32 agreement within $10^{-3}$, which the pilot's records pass.

\paragraph{Amendment 2: a bigger-model arm (during development, before any of its computation).} We
repeated the one-site arm in GPT-2 Medium on a cloud GPU (an H100), with the same prompts (the two models share
GPT-2 Small's tokeniser byte for byte), the same settings rules (candidate layers \CandMedium{}, the same relative
depths) and predictions 1--3, under a budget of \ArmBudget. DAS training used TF32 matrix products to fit the budget;
every census, margin and gradient used full float32. A first stage reproduced \NReproQueries{} saved GPT-2 Small
margins within \CheckReproGap{} (limit $10^{-3}$); refitted with TF32, the same GPT-2 Small alignment reached a best
loss of \TfLoss{} against \LocalLoss{} in float32, with \TfAgreePct\% of its single decisions equal. A cost rule fixed
in the amendment then gave GPT-2 Medium \NTasksMedium{} of the planned \NTasksConf{} confirmation tasks, and ruled out
GPT-2 Large. The arm spent \ArmSpent{} in all.

\paragraph{A development-informed addendum (before the freeze).} After reading the development results
(\cref{app:gpt2-development}), we added a reported secondary comparison of joint tests with single tests drawn
from all single queries, including those in which the swapped variable keeps its value. The hashed predictions were
not changed. The same addendum recorded that two items planned for development had not been written: an independent
float64 verifier and two simulated small-budget tests. The verifier was written after the freeze and is labelled post
hoc (\cref{app:gpt2-verification}); the simulated tests were not run.

\paragraph{Other departures from the protocol.} The learning-rate search ran at layer 9 of GPT-2 Small, the third of
the four candidates, where the protocol said layer 8; the winning rate won by a wide margin
(\cref{app:gpt2-development}). The development run was paused for a power cut, and the confirmation runs shared the
computer with other work; neither changes a computation.

\section{Study 4: development}\label{app:gpt2-development}

On the development tasks, rank-2 DAS failed none of the full-domain single queries at any candidate layer of either
model, so the tie rule chose the latest layer (\LayerSmall{} in GPT-2 Small, \LayerMedium{} in GPT-2 Medium); the
learning rate was \LrSmall{} and the number of steps \StepsSmall{} in both. In GPT-2 Small the learning-rate
search gave a mean best loss of \LrLossBest{} for the chosen rate, against \LrLossSecond{} and \LrLossThird{} for the
other two. The development results that led to the addendum are those of \cref{sec:gpt2-results}, on fewer tasks: in
GPT-2 Small no alignment passed every single query, rank-2 and rank-4 DAS failed \DevUnchangedRTwoPct\% and
\DevUnchangedRFourPct\% of the single position queries whose donor shares the recipient's position, and joint minus
full-domain single detection at \Budget{} calls was \DevContrastSmall.

\section{Study 4: verification}\label{app:gpt2-verification}

A second implementation of GPT-2 Small in NumPy float64, written after the freeze and sharing no model code with the
study, recomputed natural states from the tokens and then the margins of a sample of saved queries from the saved
subspaces: every query of every alignment that failed only jointly, every query within 0.05 of the tie band, and 2,000
random queries per phase, across arms a, a$'$, b and c. Of \VerifyQueriesConf{} confirmation queries,
\VerifyDisagreeConf{} changed decision; the largest margin difference was \VerifyMaxGapConf. GPT-2 Medium's arrays
were not re-verified; the first stage of its arm reproduced saved GPT-2 Small margins before any GPT-2 Medium
computation (\cref{app:gpt2-history}).

\section{Reproductions of published demonstrations}\label{app:reproductions}

\paragraph{Divergent representations.}
We ran the released sanity-check script for Section~4.2 of \citet{grant2026divergent} (repository commit
\texttt{f2548d2}) unchanged; all its assertions pass. In its hand-set ReLU layer the natural class-A states score
\GrantNatA{} and the class-B states score zero; adding the mean difference vector to the two class-B states raises
their scores to \GrantPatched{}, and the first of the two edits does so through a unit that is silent on every
natural input. This confirms the released example; it does not test the paper's trained models.

\paragraph{Causal abstraction metrics.}
We ran Experiment~2 (an XOR backup path) of the released code of \citet{meloux2026validating} (commit
\texttt{d91f65b}) through its own experiment factory, \CaeRuns{} runs per condition with 500 samples each, which
are the first \CaeRuns{} of the 100 official runs. The two faithfulness-aware metrics separate the valid model from
the backup-path model (two-sided Mann--Whitney $p=\CaePDown$ and \CaePUp), while the metrics without faithfulness
and interchange intervention accuracy do not ($p\ge\CaePOtherMin$), as the paper reports. The valid model scores
\CaeValidDown{} and \CaeValidUp{} on the two faithfulness-aware metrics (\CaeValidDownNf{} and \CaeValidUpNf{} without
faithfulness; interchange intervention accuracy \CaeValidIia) rather than zero. In the released script, the valid
model's mediator ($m$ in the paper, \texttt{a} in the code) does not respond to interventions: line 107 of
\texttt{systems/11\_controlled\_experiments.py} sets it to the input (\texttt{a = x.copy()}) whatever intervention is
requested, while the library passes the intervention correctly. With this line changed so that the intervention is
applied, every metric is exactly zero on the valid model, and only the faithfulness-aware metrics are non-zero on the
backup-path model (\CaeFixedLeakDown{} and \CaeFixedLeakUp). On 6 October 2026 the file in the authors' repository
was identical, byte for byte, at the cited commit and on its default branch.

\section{Records}\label{app:records}

The project folders contain, for each study, the frozen protocol and hashes, all networks, alignments, query
records and censuses, the analysis and verification scripts with their outputs, and the scripts that regenerate
every number, table and figure of this paper from those outputs. For Study~4 they also hold the prompt lists, the
development and confirmation records of both models, the cloud arm's check, cost and spending records, and the
verifier's outputs; GPT-2 Medium's full arrays stay with the cloud provider's storage. An earlier pilot with two
hand-built Boolean programs reached the same conclusion as Study~1 and is kept with its records; it is not used here.
